% Definición de homomorfismo

\begin{frame}
    \frametitle{Aplicaciones compatibles con las operaciones}

    Varias aplicaciones que ya hemos usado \alert{respetan} la suma y el
    producto por escalares:
    \begin{itemize}
        \item la aplicación de coordenadas $c_B\colon V\to k^n$ cumple
        $(u+u')_B=u_B+u'_B$ y $(\lambda u)_B=\lambda u_B$ (Tema~2);
        \item la derivación $p\mapsto p'$ de polinomios cumple
        $(p+q)'=p'+q'$ y $(\lambda p)'=\lambda p'$;
        \item la trasposición de matrices cumple $(A+B)^t=A^t+B^t$ y
        $(\lambda A)^t=\lambda A^t$.
    \end{itemize}

    \medskip

    Estas aplicaciones \alert{transportan} la estructura de espacio
    vectorial de un espacio a otro. Son el objeto de estudio de este
    tema.
\end{frame}


\begin{frame}
    \frametitle{Homomorfismos de espacios vectoriales}

    \begin{definition}
        Sean $V$ y $V'$ espacios vectoriales sobre $k$. Un
        \alert{homomorfismo} (o \alert{aplicación lineal}) $f\colon V\to V'$ de $V$ en $V'$
        es una aplicación tal que, para todo
        $u,u'\in V$ y todo $\lambda\in k$,
        \begin{enumerate}[(1)]
            \item $f(u+u')=f(u)+f(u')$;
            \item $f(\lambda u)=\lambda f(u)$.
        \end{enumerate}
        Denotamos por $\Hom(V,V')$ el conjunto de los homomorfismos de $V$
        en $V'$.
    \end{definition}

    \medskip

    En (1) y (2), la suma y el producto por escalares del primer miembro
    son los de $V$, y los del segundo miembro son los de $V'$.
\end{frame}


\begin{frame}
    \frametitle{Aplicaciones definidas por matrices}

    \begin{proposition}
        Sea $A\in M_{m\times n}(k)$. La aplicación
        $f_A\colon k^n\to k^m$, $f_A(u)=Au$, es un homomorfismo.
    \end{proposition}

    \vspace{-6pt}

    \begin{proof}
        Por las propiedades del producto de matrices (Tema~1), para todo
        $u,u'\in k^n$ y $\lambda\in k$,
        \[
            A(u+u')=Au+Au',
            \qquad
            A(\lambda u)=\lambda(Au).\qedhere
        \]
    \end{proof}

    \vspace{-6pt}

    Veremos que \alert{todos} los homomorfismos entre espacios de
    dimensión finita son, en esencia, de este tipo.
\end{frame}


\begin{frame}
    \frametitle{Ejemplo: un giro en el plano}

    \begin{columns}[T]
        \begin{column}{0.58\textwidth}
            \begin{example}
                Sea $\theta\in\R$ y
                \[
                    R_\theta=\begin{pmatrix}
                        \cos\theta&-\sin\theta\\
                        \sin\theta&\cos\theta
                    \end{pmatrix}.
                \]
                La aplicación $f_\theta=f_{R_\theta}\colon\R^2\to\R^2$,
                $f_\theta(u)=R_\theta u$, es el \alert{giro} de ángulo $\theta$.
            \end{example}
        \end{column}
        \begin{column}{0.38\textwidth}
            \centering
            \begin{tikzpicture}[scale=1.5]
                \draw[->] (-0.3,0) -- (1.35,0);
                \draw[->] (0,-0.3) -- (0,1.35);
                \draw[gray,dashed] (1,0) arc (0:90:1);
                \draw[->,very thick] (0,0) -- (1,0)
                    node[below right] {\footnotesize $e_1$};
                \draw[->,very thick] (0,0) -- (0,1)
                    node[above left] {\footnotesize $e_2$};
                \draw[->,very thick,alertcolor] (0,0) -- (0.866,0.5)
                    node[right] {\footnotesize $f_\theta(e_1)$};
                \draw[->,very thick,alertcolor] (0,0) -- (-0.5,0.866)
                    node[above left] {\footnotesize $f_\theta(e_2)$};
                \draw (0.4,0) arc (0:30:0.4);
                \node at (0.52,0.12) {\scriptsize $\theta$};
            \end{tikzpicture}
        \end{column}
    \end{columns}

    \medskip

    Por la proposición anterior, $f_\theta$ es un homomorfismo. Sus columnas
    son
    \[
        f_\theta(e_1)=(\cos\theta,\sin\theta),
        \qquad
        f_\theta(e_2)=(-\sin\theta,\cos\theta),
    \]
    los vectores $e_1$ y $e_2$ girados un ángulo $\theta$ en sentido
    antihorario. Como $f_\theta$ preserva combinaciones lineales, queda
    determinada por estas dos imágenes.
\end{frame}


\begin{frame}
    \frametitle{Ejemplos básicos}

    \begin{example}
        Sean $V$ y $V'$ espacios vectoriales sobre $k$.
        \begin{enumerate}
            \item La aplicación \alert{nula} $0\colon V\to V'$,
            $u\mapsto 0$.
            \item La aplicación \alert{identidad}
            $\id_V\colon V\to V$, $u\mapsto u$.
            \item La \alert{homotecia} de razón $\lambda\in k$,
            $V\to V$, $u\mapsto\lambda u$.
            \item La \alert{inclusión} $U\hookrightarrow V$ de un
            subespacio $U$ de $V$.
            \item La aplicación de coordenadas
            $c_B\colon V\to k^n$ respecto de una base $B$ de $V$.
        \end{enumerate}
    \end{example}

    La linealidad de $c_B$ es una proposición del Tema~2; la de los demás
    ejemplos es inmediata.
\end{frame}


\begin{frame}
    \frametitle{Proyecciones, inclusiones y proyección canónica}

    \begin{example}
        \begin{enumerate}
            \item Las \alert{proyecciones}
            $V\times V'\to V$, $(u,w)\mapsto u$, y
            $V\times V'\to V'$, $(u,w)\mapsto w$.
            \item Las \alert{inclusiones}
            $V\to V\times V'$, $u\mapsto(u,0)$, y
            $V'\to V\times V'$, $w\mapsto(0,w)$.
            \item La \alert{proyección canónica}
            $\pi\colon V\to V/U$, $\pi(u)=u+U$, sobre el cociente de $V$
            por un subespacio $U$.
        \end{enumerate}
    \end{example}

    La linealidad de $\pi$ es la definición de las operaciones en $V/U$:
    \begin{align*}
        \pi(u+u')&=(u+u')+U=(u+U)+(u'+U)=\pi(u)+\pi(u'),
        \\
        \pi(\lambda u)&=(\lambda u)+U=\lambda(u+U)=\lambda\pi(u).
    \end{align*}
\end{frame}


\begin{frame}
    \frametitle{Ejemplos en polinomios, matrices y funciones}

    \begin{example}
        \begin{enumerate}
            \item La \alert{derivación}
            $D\colon k[x]\to k[x]$, $D(a_0+a_1x+\cdots+a_nx^n)
            =a_1+2a_2x+\cdots+na_nx^{n-1}$.
            \item La \alert{evaluación} en $a\in k$,
            $\operatorname{ev}_a\colon k[x]\to k$, $p\mapsto p(a)$.
            \item La \alert{multiplicación} por $x$,
            $k[x]\to k[x]$, $p\mapsto xp$.
            \item La \alert{traza} $\operatorname{tr}\colon M_n(k)\to k$,
            $(a_{ij})\mapsto a_{11}+\cdots+a_{nn}$.
            \item La \alert{trasposición}
            $M_{m\times n}(k)\to M_{n\times m}(k)$, $A\mapsto A^t$.
            \item La \alert{multiplicación} por una matriz fija $A\in M_n(k)$,
            $M_n(k)\to M_n(k)$, $X\mapsto AX$, y el \alert{conmutador}
            $X\mapsto AX-XA$.
            \item La \alert{evaluación} en un punto $x_0\in X$,
            $k^X\to k$, $f\mapsto f(x_0)$.
        \end{enumerate}
    \end{example}

    Todas se comprueban con las propiedades de la suma y del producto por
    escalares en cada espacio.
\end{frame}


\begin{frame}
    \frametitle{Una caracterización con un solo requisito}

    \begin{proposition}
        Una aplicación $f\colon V\to V'$ es un homomorfismo si y solo si
        para todo $u,u'\in V$ y $\lambda,\mu\in k$,
        \[
            f(\lambda u+\mu u')=\lambda f(u)+\mu f(u').
        \]
    \end{proposition}

    \begin{proof}
        Si $f$ es un homomorfismo,
        $f(\lambda u+\mu u')=f(\lambda u)+f(\mu u')=\lambda f(u)+\mu f(u')$.
        Recíprocamente, tomando $\lambda=\mu=1$ se obtiene la propiedad (1) de la definición y tomando
        $\mu=0$ se obtiene (2).
    \end{proof}
\end{frame}


\begin{frame}
    \frametitle{Los homomorfismos preservan combinaciones lineales}

    \begin{proposition}
        Sea $f\colon V\to V'$ un homomorfismo. Para todo
        $u_1,\ldots,u_r\in V$ y $\lambda_1,\ldots,\lambda_r\in k$,
        \[
            f(\lambda_1u_1+\cdots+\lambda_ru_r)
            =\lambda_1f(u_1)+\cdots+\lambda_rf(u_r).
        \]
        Además, $f(0)=0$ y $f(-u)=-f(u)$ para todo $u\in V$.
    \end{proposition}

    \begin{proof}
        La fórmula se obtiene por inducción sobre $r$: para $r=1$ es (2) y,
        si $r\geq2$, aplicando (1), la hipótesis de inducción y (2),
        \begin{align*}
            f(\lambda_1u_1+\cdots+\lambda_ru_r)
            &=f\bigl((\lambda_1u_1+\cdots+\lambda_{r-1}u_{r-1})+\lambda_ru_r\bigr)\\
            &=f(\lambda_1u_1+\cdots+\lambda_{r-1}u_{r-1})+f(\lambda_ru_r)\\
            &=\lambda_1f(u_1)+\cdots+\lambda_{r-1}f(u_{r-1})+\lambda_rf(u_r).
        \end{align*}
        Además, $f(0)=f(0\cdot 0)=0\cdot f(0)=0$ y
        $f(-u)=f((-1)u)=(-1)f(u)=-f(u)$.
    \end{proof}
\end{frame}


\begin{frame}
    \frametitle{Contraejemplo: las traslaciones}

    \begin{example}
        Sea $a\in V$ no nulo. La \alert{traslación}
        $f\colon V\to V$, $f(u)=u+a$, no es un homomorfismo, porque
        \[
            f(0)=a\neq0.
        \]
    \end{example}

    \medskip

    Así, $f(0)=0$ es una condición \alert{necesaria}, pero no suficiente.

\end{frame}


\begin{frame}
    \frametitle{Contraejemplos: se cumple una condición, pero no la otra}

    \begin{example}
        \begin{enumerate}
            \item $f\colon\R\to\R$, $f(u)=u^2$ cumple $f(0)=0$, pero
            $f(1+1)=4\neq 2=f(1)+f(1)$.
            \item $f\colon\R^2\to\R$,
            $f(u_1,u_2)=\sqrt[3]{u_1^3+u_2^3}$ cumple
            $f(\lambda u)=\lambda f(u)$, pero no (1):
            \[
                f(1,0)+f(0,1)=2\neq\sqrt[3]{2}=f(1,1).
            \]
            \item $f\colon\C\to\C$, $f(z)=\bar z$ cumple
            $f(z+z')=f(z)+f(z')$, pero no (2) si $k=\C$:
            $f(iz)=-i\bar z\neq i\bar z=i f(z)$ si $z\neq0$.
            (Sí es un homomorfismo si consideramos $\C$ como espacio
            vectorial sobre $\R$.)
            \item $\det\colon M_2(\R)\to\R$ no es un homomorfismo:
            $\det(I+I)=4\neq 2=\det I+\det I$.
        \end{enumerate}
    \end{example}

    Las dos condiciones de la definición de homomorfismo son \alert{independientes}.
\end{frame}


\begin{frame}
    \frametitle{La composición de homomorfismos}

    \begin{proposition}
        Sean $f\colon V\to V'$ y $g\colon V'\to V''$ homomorfismos. Entonces
        $g\circ f\colon V\to V''$ es un homomorfismo.
    \end{proposition}

    \begin{proof}
        Para todo $u,u'\in V$ y $\lambda\in k$,
        \[
            (g\circ f)(u+u')=g\bigl(f(u)+f(u')\bigr)
            =(g\circ f)(u)+(g\circ f)(u'),
        \]
        \[
            (g\circ f)(\lambda u)=g\bigl(\lambda f(u)\bigr)
            =\lambda(g\circ f)(u).\qedhere
        \]
    \end{proof}
\end{frame}


\begin{frame}
    \frametitle{La inversa de un homomorfismo biyectivo}

    \begin{proposition}
        Sea $f\colon V\to V'$ un homomorfismo biyectivo. Entonces su inversa
        $f^{-1}\colon V'\to V$ es un homomorfismo.
    \end{proposition}

    \begin{proof}
        Sean $w,w'\in V'$ y $\lambda\in k$. Como $f$ es un
        homomorfismo,
        \begin{align*}
            f\bigl(f^{-1}(w)+f^{-1}(w')\bigr)
            &=f\bigl(f^{-1}(w)\bigr)+f\bigl(f^{-1}(w')\bigr)=w+w',\\
            f\bigl(\lambda f^{-1}(w)\bigr)
            &=\lambda f\bigl(f^{-1}(w)\bigr)=\lambda w.
        \end{align*}
        Aplicando $f^{-1}$ a ambas igualdades,
        \begin{align*}
            f^{-1}(w)+f^{-1}(w')&=f^{-1}(w+w'),\\
            \lambda f^{-1}(w)&=f^{-1}(\lambda w).\qedhere
        \end{align*}
    \end{proof}
\end{frame}


\begin{frame}
    \frametitle{Isomorfismos, endomorfismos y automorfismos}

    \begin{definition}
        Sea $f\colon V\to V'$ un homomorfismo.
        \begin{itemize}
            \item $f$ es un \alert{monomorfismo} si es inyectivo. Lo
            indicamos $f\colon V\hookrightarrow V'$.
            \item $f$ es un \alert{epimorfismo} si es sobreyectivo. Lo
            indicamos $f\colon V\twoheadrightarrow V'$.
            \item $f$ es un \alert{isomorfismo} si es biyectivo, es decir,
            si es monomorfismo y epimorfismo. Lo indicamos
            $f\colon V\xrightarrow{\,\sim\,}V'$. 
            \item $f$ es un \alert{endomorfismo} si $V'=V$. Escribimos
            $\End(V)=\Hom(V,V)$.
            \item $f$ es un \alert{automorfismo} si es un endomorfismo y
            un isomorfismo.
        \end{itemize}
    \end{definition}

    \medskip

    Si existe
    un isomorfismo $V\to V'$, decimos que $V$ y $V'$ son
    \alert{isomorfos} y escribimos $V\cong V'$.
    Por la proposición anterior, la identidad es un isomorfismo, la
    inversa de un isomorfismo es un isomorfismo y la composición de
    isomorfismos es un isomorfismo. Por tanto, \alert{ser isomorfos} es
    una relación de equivalencia entre espacios vectoriales.

\end{frame}


\begin{frame}
    \frametitle{Ejemplos de isomorfismos}

    \begin{example}
        \begin{enumerate}
            \item Si $B=\{v_1,\ldots,v_n\}$ es una base de $V$ con $n$ vectores, la
            aplicación de coordenadas $c_B\colon V\to k^n$ es un
            isomorfismo (Tema~2). Su inversa es
            $(\lambda_1,\ldots,\lambda_n)\mapsto\lambda_1v_1+\cdots+\lambda_nv_n$.
            \item $k_n[x]\cong k^{n+1}$, mediante
            $a_0+a_1x+\cdots+a_nx^n\mapsto(a_0,a_1,\ldots,a_n)$.
            \item $M_{m\times n}(k)\cong k^{mn}$, escribiendo las
            entradas de cada matriz en una única fila.
            \item $M_2(\R)\to M_2(\R)$, $A\mapsto A^t$, es un automorfismo,
            y es su propia inversa pues $(A^t)^t=A$.
            \item $k^m\times k^n\cong k^{m+n}$, mediante la concatenación de coordenadas
            $((u_1,\ldots,u_m),(u'_1,\ldots,u'_n))\mapsto(u_1,\ldots,u_m,u'_1,\ldots,u'_n)$.
        \end{enumerate}
    \end{example}
\end{frame}

\begin{frame}
    \frametitle{Homomorfismos que coinciden en un sistema generador}

    \begin{proposition}
        Dos homomorfismos $f,g\colon V\to V'$ que coinciden en un sistema
        generador de $V$ son iguales.
    \end{proposition}

    \begin{proof}
        Sean $v_1,\ldots,v_r$ generadores de $V$ con $f(v_j)=g(v_j)$ para
        todo $j$, y sea $u\in V$. Existen $\lambda_1,\ldots,\lambda_r\in k$
        con $u=\lambda_1v_1+\cdots+\lambda_rv_r$. Como $f$ y $g$ preservan
        combinaciones lineales,
        \begin{align*}
            f(u)&=f(\lambda_1v_1+\cdots+\lambda_rv_r)\\
            &=\lambda_1f(v_1)+\cdots+\lambda_rf(v_r)\\
            &=\lambda_1g(v_1)+\cdots+\lambda_rg(v_r)\\
            &=g(\lambda_1v_1+\cdots+\lambda_rv_r)\\
            &=g(u).
        \end{align*}
        Por tanto, $f(u)=g(u)$ para todo $u\in V$, es decir, $f=g$.
    \end{proof}
\end{frame}

\begin{frame}[allowframebreaks, t]
    \frametitle{Un homomorfismo se determina por la imagen de una base}

    \begin{theorem}
        Sea $B=\{v_1,\ldots,v_n\}$ una base de $V$ y sean
        $w_1,\ldots,w_n\in V'$ \alert{cualesquiera}. Existe un único
        homomorfismo $f\colon V\to V'$ tal que
        \[
            f(v_j)=w_j,\qquad j=1,\ldots,n.
        \]
        Si $u_B=(\alpha_1,\ldots,\alpha_n)$, entonces
        $f(u)=\alpha_1w_1+\cdots+\alpha_nw_n$.
    \end{theorem}

        Al elegir las imágenes $w_1,\ldots,w_n$ de una base \alert{no hay
    ninguna restricción}: pueden ser nulas, repetirse o ser linealmente
    dependientes.


    \framebreak

    \begin{proof}
        \alert{Unicidad.} Si $f$ cumple $f(v_j)=w_j$ y
        $u=\alpha_1v_1+\cdots+\alpha_nv_n$, como $f$ preserva
        combinaciones lineales,
        $f(u)=\alpha_1w_1+\cdots+\alpha_nw_n$.

        \medskip

        \alert{Existencia.} Si $u_B=(\alpha_1,\ldots,\alpha_n)$,
        definimos
        \[
            f(u)=\alpha_1w_1+\cdots+\alpha_nw_n.
        \]
        Está bien definida porque las coordenadas de $u$ respecto de $B$ son
        únicas.

        \medskip

        \alert{$f(v_j)=w_j$.} Como $(v_j)_B=e_j$, es decir, todas las
        coordenadas de $v_j$ son $0$ salvo la $j$-ésima, que es $1$,
        \[
            f(v_j)=0\cdot w_1+\cdots+1\cdot w_j+\cdots+0\cdot w_n=w_j.
        \]

        \framebreak

        \alert{$f$ es un homomorfismo.} Sean $u,u'\in V$ y $\lambda\in k$,
        con
        \[
            u_B=(\alpha_1,\ldots,\alpha_n),\qquad u'_B=(\beta_1,\ldots,\beta_n).
        \]
        Como las coordenadas respecto de $B$ son aditivas y homogéneas,
        \[
            (u+u')_B=(\alpha_1+\beta_1,\ldots,\alpha_n+\beta_n),
            \qquad
            (\lambda u)_B=(\lambda \alpha_1,\ldots,\lambda \alpha_n).
        \]
        Por la definición de $f$,
        \begin{align*}
            f(u+u')&=(\alpha_1+\beta_1)w_1+\cdots+(\alpha_n+\beta_n)w_n\\
            &=(\alpha_1w_1+\cdots+\alpha_nw_n)+(\beta_1w_1+\cdots+\beta_nw_n)\\
            &=f(u)+f(u'),\\[1ex]
            f(\lambda u)&=\lambda \alpha_1w_1+\cdots+\lambda \alpha_nw_n\\
            &=\lambda(\alpha_1w_1+\cdots+\alpha_nw_n)\\
            &=\lambda f(u).\qedhere
        \end{align*}
    \end{proof}
\end{frame}


\begin{frame}
    \frametitle{Ejemplo: homomorfismo dado por la imagen de una base}

    \begin{example}
        Sea $B=\{(1,1),(1,-1)\}\subset\R^2$, que es una base pues
        \[
            \det \begin{pmatrix}1&1\\1&-1\end{pmatrix}=-2\neq0.
        \]
        Hallemos el
        homomorfismo $f\colon\R^2\to\R^3$ con
        \[
            f(1,1)=(1,0,2),\qquad f(1,-1)=(0,1,1).
        \]
        Si $C$ es la base canónica, $M(B,C)$ es la matriz anterior y 
        \[M(C,B)=M(B,C)^{-1}=\frac{1}{2}\begin{pmatrix}1&1\\1&-1\end{pmatrix},\]
        que hemos calculado como la traspuesta de la adjunta de partida por su determinante.
    \end{example}

\end{frame}

\begin{frame}
    \frametitle{Ejemplo: homomorfismo dado por la imagen de una base}
    
    Así, las coordenadas de $u=(u_1,u_2)$ en $B$ son
    \[
        u_B=M(C,B)\begin{pmatrix}u_1\\u_2\end{pmatrix}
        =\begin{pmatrix}\frac{u_1+u_2}{2}\\[0.5ex]\frac{u_1-u_2}{2}\end{pmatrix},
    \]
    así que
    \[
        f(u_1,u_2)=\tfrac{u_1+u_2}{2}(1,0,2)+\tfrac{u_1-u_2}{2}(0,1,1)
        =\Bigl(\tfrac{u_1+u_2}{2},\ \tfrac{u_1-u_2}{2},\ \tfrac{3u_1+u_2}{2}\Bigr).
    \]

    \medskip

    En particular, como $(3,1)=2(1,1)+1\cdot(1,-1)$,
    \[
        f(3,1)=2(1,0,2)+(0,1,1)=(2,1,5).
    \]
\end{frame}


\begin{frame}
    \frametitle{Ejemplo: homomorfismos desde polinomios}

    \begin{example}
        Existe un único homomorfismo $f\colon\R_2[x]\to\R^2$ con
        \[
            f(1)=(1,0),\qquad f(x)=(1,1),\qquad f(x^2)=(0,2).
        \]
    \end{example}

    Como las coordenadas de $a+bx+cx^2$ en $\{1,x,x^2\}$ son $(a,b,c)$,
    \[
        f(a+bx+cx^2)=a(1,0)+b(1,1)+c(0,2)=(a+b,\ b+2c).
    \]

    \medskip

    Para $p(x)=2-x+3x^2$ se obtiene $f(p)=(1,5)$.

    \medskip

    En este ejemplo, los tres vectores $(1,0),(1,1),(0,2)$ de $\R^2$ son
    linealmente dependientes, pero $f$ existe de todos modos: lo que
    importa es que $1,x,x^2$ sean una \alert{base} de $\R_2[x]$.
\end{frame}
